\section{Unified Evaluation Harness}
\label{sec:harness}

Our first contribution is a reproducible evaluation harness for comparing agent conditions on equal footing. The harness is designed around three principles: \textbf{pinned versions} (every benchmark is locked to a specific git commit or release), \textbf{unified logging} (every run produces a manifest with model metadata, config, and cost accounting), and \textbf{within-benchmark paired comparison} (all conditions are evaluated on the same task instances, enabling paired statistical analysis).

% ------------------------------------------------------------
%  3.1  Benchmark Suite
% ------------------------------------------------------------

\subsection{Benchmark Suite}
\label{sec:harness:benchmarks}

We evaluate on two benchmarks that span diverse agent task profiles:

\begin{itemize}[leftmargin=1.5em, itemsep=1pt, parsep=0pt]

  \item \textbf{AgentBench-OpenClaw} (AB-OC): A general multi-step tool-use benchmark with 40 tasks across 7 domains (file creation, research, data analysis, multi-step, memory, error handling, tool efficiency). Tasks average 6--10 steps with multiple tool calls. We use the pinned OpenClaw skill release with the standard task split.

  \item \textbf{PinchBench} (PB): A benchmark of messy real-world productivity tasks including calendar management, email triage, market research, code generation, and cross-session memory. Tasks involve moderate tool use (2--5 calls) but heavy state management requirements. We use the pinned \texttt{v0.9} release with all 23 tasks.

\end{itemize}

Table~\ref{tab:benchmark-suite} in \S\ref{sec:exp:setup} reports the full suite metadata. All benchmarks provide deterministic gold answers or structured grading scripts; we use strict equality for factual tasks and LLM-as-judge (GPT-4o) for synthesis tasks, following established protocols~\citep{zheng2023judging}.

% ------------------------------------------------------------
%  3.2  Agent Conditions
% ------------------------------------------------------------

\subsection{Agent Conditions}
\label{sec:harness:conditions}

We compare the following conditions, all evaluated on the nanobot OpenClaw agent runtime with doubao-seed-1.8 as the base model:

\begin{enumerate}[leftmargin=2em, label=(\alph*), itemsep=0pt, parsep=0pt]

  \item \textbf{Base}: The unmodified OpenClaw agent with greedy decoding, temperature 0.0, and max 40 steps.

  \item \textbf{+Memory}: Base agent augmented with a task-local episodic memory buffer that records intermediate facts, file paths, and verified conclusions. Memory is written after each tool call and read before each new tool invocation. \emph{[Planned; not yet experimentally evaluated]}

  \item \textbf{+Control}: Base agent with plan-first prompting (explicitly generate a step plan before acting), bounded replanning (replan every 5 steps), and failure-triggered reflection (if a tool returns an error, generate an explicit diagnosis before retrying). \emph{[Planned; not yet experimentally evaluated]}

  \item \textbf{+Collab}: Base agent with a minimal two-role structure: an \emph{executor} that calls tools and a \emph{verifier} that checks intermediate conclusions against the task goal. Roles communicate via a shared message buffer and the verifier can request a rollback to a previous checkpoint. \emph{[Planned; not yet experimentally evaluated]}

  \item \textbf{+ProcSupport}: Base agent with on-demand access to structured procedural support cards for each benchmark domain. Cards contain step-by-step domain procedures (e.g., ``how to conduct a PubMed search for clinical trials'') and are retrieved via semantic similarity at task start. \emph{[Planned; not yet experimentally evaluated]}

  \item \textbf{+BestTrainFree}: The best combination of training-free recipes as determined on preliminary analysis (memory + control). \emph{[Planned; not yet experimentally evaluated]}

  \item \textbf{SFT-100}, \textbf{SFT-300}, \textbf{SFT-1000}: LoRA-SFT variants trained on 100, 300, and 1,000 corrected trajectories respectively, sampled from failures of the Base condition. Training uses LoRA rank 16, learning rate $2\times10^{-4}$, batch size 8, 2 epochs. \emph{[Planned; not yet experimentally evaluated]}

\end{enumerate}

% ------------------------------------------------------------
%  3.3  Cost Accounting and Metrics
% ------------------------------------------------------------

\subsection{Cost Accounting and Metrics}
\label{sec:harness:metrics}

Every run records: total prompt tokens, completion tokens, tool calls, steps, wall-clock time, and estimated API cost (USD). Our primary metric is \textbf{within-benchmark paired delta}:

\begin{align}
  \Delta_b(c) = \frac{1}{|\mathcal{T}_b|}\sum_{t\in\mathcal{T}_b}\bigl[\mathrm{score}_c(t) - \mathrm{score}_{\mathrm{Base}}(t)\bigr],
  \label{eq:paired-delta}
\end{align}

where $\mathcal{T}_b$ is the task set for benchmark $b$, and $\mathrm{score} \in \{0, 1\}$ is binary correctness. We report mean and 95\% bootstrap confidence intervals over 1,000 resamples.

We also track \textbf{negative interference rate} (NIR): the fraction of tasks where a repair condition scores strictly below Base:

\begin{align}
  \mathrm{NIR}(c) = \frac{1}{|\mathcal{T}_b|}\sum_{t\in\mathcal{T}_b} \mathbf{1}\bigl[\mathrm{score}_c(t) < \mathrm{score}_{\mathrm{Base}}(t)\bigr].
  \label{eq:nir}
\end{align}

And \textbf{category repair rate} (CRR) for each failure category $k \in \{A, B, C, D, E\}$:

\begin{align}
  \mathrm{CRR}_k(c) = P\bigl(\mathrm{score}_c(t) = 1 \mid \mathrm{score}_{\mathrm{Base}}(t) = 0, \\
  \qquad\qquad\qquad\;\; \mathrm{cat}(t) = k\bigr).
  \label{eq:crr}
\end{align}

\begin{definit}[Failure Category]\label{def:failure-category}
For each task run where $\mathrm{score}_{\mathrm{Base}}(t) = 0$, the \emph{dominant failure category} $\mathrm{cat}(t) \in \{A, B, C, D, E\}$ is the single category that best explains why the agent failed, as determined by structured annotation (\S\ref{sec:taxonomy:annotation}).
\end{definit}
